% GENERATED FILE. Edit canonical lessons, then run npm run generate:books.
% canonical-source-hash: fnv1a64:cad14e31ca8e4ab1
% canonical-lessons: ZH-C120-waimai, ZH-C120-ding, ZH-C120-dizhi, ZH-C120-dabao, ZH-C120-dianxin

\chapter{Takeaway}
\label{ch:zh-takeaway}
\hlchaptermodality{120}

\begin{chapteropening}
\textbf{By the end of this chapter:} \emph{I can order a takeaway, give my address, ask to take food home, and ask for dim sum.}

Complete the last lesson of chapter 120: I can order a takeaway, give my address, ask to take food home, and ask for dim sum.
\end{chapteropening}

\section[wàimài]{wàimài — takeaway; food delivered}

\label{lesson:ZH-C120-waimai}



\begin{quote}
Before the new one: say the Mandarin for carrier bag, then the Mandarin for luggage.
\end{quote}

\subsection*{You'll want to know: wàimài}
\emph{wàimài} — \textquotedblleft{}takeaway; food delivered\textquotedblright{}.

\begin{quote}
\emph{wǒmen jīntiān wǎnshang chī wàimài} — We're having a takeaway tonight.
\end{quote}

Its characters come later; for now it is a word for the ear.

\begin{cousinweb}
\emph{wàimài} is built from \zh{外} 'outside' and mài 'to sell', both inherited from Middle Chinese, \emph{ngwajH} and \emph{meaH} (Guangyun, in Baxter's transcription): food sold to be eaten somewhere else.
\end{cousinweb}

\begin{grammarlens}[title={What you've built}]
One more word for asking for what I want.
\end{grammarlens}

\subsection*{Your turn}
\begin{itemize}
\raggedright
  \item \emph{Say it:} \emph{wàimài}
  \item \emph{Say it:} \emph{wàimài}, once more
  \item \emph{Say it:} \emph{wǒmen jīntiān wǎnshang chī wàimài} — We're having a takeaway tonight.
  \item \emph{Recall:} say the Mandarin for carrier bag, then the Mandarin for luggage, then say \emph{wàimài} again
\end{itemize}

\begin{tcolorbox}[breakable,colback=teal!4,colframe=teal!35!black,title={Before you move on}]
What does wàimài mean? (\textquotedblleft{}Takeaway; food delivered\textquotedblright{}.) Say it once more.
\end{tcolorbox}
\section[dìng]{dìng — to make a booking; to order}

\label{lesson:ZH-C120-ding}



\begin{quote}
Before the new one: say the Mandarin for luggage, then the Mandarin for takeaway.
\end{quote}

\subsection*{You'll want to know: dìng}
\emph{dìng} — \textquotedblleft{}to make a booking; to order\textquotedblright{}.

\begin{quote}
\emph{wǒ xiǎng dìng yí ge fángjiān} — I'd like to book a room.
\end{quote}

Its characters come later; for now it is a word for the ear.

\begin{cousinweb}
\emph{dìng} is inherited from Middle Chinese \emph{dengX} (Guangyun, in Baxter's transcription), glossed 'to discuss and settle', which regularly gives dìng. Booking is the sense today, probably by way of settling an arrangement, but the step is not given there.
\end{cousinweb}

\begin{grammarlens}[title={What you've built}]
One more word for asking for what I want.
\end{grammarlens}

\subsection*{Your turn}
\begin{itemize}
\raggedright
  \item \emph{Say it:} \emph{dìng}
  \item \emph{Say it:} \emph{dìng}, once more
  \item \emph{Say it:} \emph{wǒ xiǎng dìng yí ge fángjiān} — I'd like to book a room.
  \item \emph{Recall:} say the Mandarin for luggage, then the Mandarin for takeaway, then say \emph{dìng} again
\end{itemize}

\begin{tcolorbox}[breakable,colback=teal!4,colframe=teal!35!black,title={Before you move on}]
What does dìng mean? (\textquotedblleft{}To make a booking; to order\textquotedblright{}.) Say it once more.
\end{tcolorbox}
\section[dìzhǐ]{dìzhǐ — address}

\label{lesson:ZH-C120-dizhi}



\begin{quote}
Before the new one: say the Mandarin for takeaway, then the Mandarin for to make a booking.
\end{quote}

\subsection*{You'll want to know: dìzhǐ}
\emph{dìzhǐ} — \textquotedblleft{}address\textquotedblright{}.

\begin{quote}
\emph{nǐ jiā de dìzhǐ shì shénme} — What's your home address?
\end{quote}

Its characters come later; for now it is a word for the ear.

\begin{cousinweb}
\emph{dìzhǐ} is built from dì 'ground' and zhǐ 'site', both inherited from Middle Chinese, \emph{dijH} and \emph{tsyiX} (Guangyun, in Baxter's transcription); the dictionary glosses the second as the foundation of a building.
\end{cousinweb}

\begin{grammarlens}[title={What you've built}]
One more word for asking for what I want.
\end{grammarlens}

\subsection*{Your turn}
\begin{itemize}
\raggedright
  \item \emph{Say it:} \emph{dìzhǐ}
  \item \emph{Say it:} \emph{dìzhǐ}, once more
  \item \emph{Say it:} \emph{nǐ jiā de dìzhǐ shì shénme} — What's your home address?
  \item \emph{Recall:} say the Mandarin for takeaway, then the Mandarin for to make a booking, then say \emph{dìzhǐ} again
\end{itemize}

\begin{tcolorbox}[breakable,colback=teal!4,colframe=teal!35!black,title={Before you move on}]
What does dìzhǐ mean? (\textquotedblleft{}Address\textquotedblright{}.) Say it once more.
\end{tcolorbox}
\section[\texorpdfstring{\zh{打包} (dǎbāo)}{打包 (dǎbāo)}]{\zh{打包} (dǎbāo) — to pack food to take away}

\label{lesson:ZH-C120-dabao}



\begin{quote}
Before the new one: say the Mandarin for to make a booking, then the Mandarin for address.
\end{quote}

\subsection*{You'll want to know: \zh{打包}}
\textbf{\zh{打包}} — \emph{dǎbāo} — \textquotedblleft{}to pack food to take away\textquotedblright{}.

\begin{quote}
\textbf{\zh{我想打包}} \emph{wǒ xiǎng dǎbāo} — Could I take this away?
\end{quote}

\begin{cousinweb}
\textbf{\zh{打包}} (\emph{dǎbāo}) is built from \zh{打} 'to strike, to do' and \zh{包} 'to wrap', both inherited from Middle Chinese, \emph{taengX} and \emph{paew} (Guangyun, in Baxter's transcription). The dictionary glosses the first 'to strike'; dǎ is not the sound \emph{taengX} would regularly give, and its everyday sense of doing is not given there.
\end{cousinweb}

\begin{grammarlens}[title={What you've built}]
One more word for asking for what I want.
\end{grammarlens}

\subsection*{Your turn}
\begin{itemize}
\raggedright
  \item \emph{Say it:} \emph{dǎbāo}
  \item \emph{Say it:} \emph{dǎbāo}, once more
  \item \emph{Say it:} \emph{wǒ xiǎng dǎbāo} — Could I take this away?
  \item \emph{Recall:} say the Mandarin for to make a booking, then the Mandarin for address, then say \emph{dǎbāo} again
\end{itemize}

\begin{tcolorbox}[breakable,colback=teal!4,colframe=teal!35!black,title={Before you move on}]
What does \zh{打包} mean? (\textquotedblleft{}To pack food to take away\textquotedblright{}.) Say it once more.
\end{tcolorbox}
\section[\texorpdfstring{\zh{点心} (diǎnxin)}{点心 (diǎnxin)}]{\zh{点心} (diǎnxin) — dim sum; light snacks and pastries}

\label{lesson:ZH-C120-dianxin}



\begin{quote}
Before the new one: say the Mandarin for address, then the Mandarin for to pack food to take away.
\end{quote}

\subsection*{You'll want to know: \zh{点心}}
\textbf{\zh{点心}} — \emph{diǎnxin} — \textquotedblleft{}dim sum; light snacks and pastries\textquotedblright{}.

\begin{quote}
\textbf{\zh{我们去外边吃点心}} \emph{wǒmen qù wàibian chī diǎnxin} — We're going out for dim sum.
\end{quote}

\begin{cousinweb}
\textbf{\zh{点心}} (\emph{diǎnxin}) is built from \zh{点} 'a dot, a little' and \zh{心} 'heart', both inherited from Middle Chinese, \emph{temX} and \emph{sim} (Guangyun, in Baxter's transcription); the dictionary glosses the first as a dot made with the brush. How the pair came to mean a light snack is not settled.
\end{cousinweb}

\begin{grammarlens}[title={What you've built}]
One more word for asking for what I want.
\end{grammarlens}

\subsection*{Your turn}
\begin{itemize}
\raggedright
  \item \emph{Say it:} \emph{diǎnxin}
  \item \emph{Say it:} \emph{diǎnxin}, once more
  \item \emph{Say it:} \emph{wǒmen qù wàibian chī diǎnxin} — We're going out for dim sum.
  \item \emph{Recall:} say the Mandarin for address, then the Mandarin for to pack food to take away, then say \emph{diǎnxin} again
\end{itemize}

\begin{tcolorbox}[breakable,colback=teal!4,colframe=teal!35!black,title={Before you move on}]
What does \zh{点心} mean? (\textquotedblleft{}Dim sum; light snacks and pastries\textquotedblright{}.) Say it once more.
\end{tcolorbox}
